\section{The identity without contact}\label{revision:offcontact}
The main proof uses the contact identity. This stronger extension
applies to every compactly supported standardized law and also provides
an identity for checking numerical implementations on nonextremal laws. Retain the definitions of $a,Q,h,\Delta$,
$J,J_1,J_2,L$, and $\mathcal R_c$ from the main text, for arbitrary
$H\in\R$ and $c>0$. Introduce the contact residuals
\begin{equation}\label{v9:core:residuals}
 \xi=\E[X\Delta(X)],\qquad
 R_{\rm supp}=\int_K\tau(t)\Delta(t)\,dt
       +\sum_{(u,v)}\tau_{u,v}(v-u)\Delta(v).
\end{equation}
The first integral is with respect to Lebesgue measure on the closed
support $K$, not to $P$. Since $\Delta$ is bounded on the convex hull
and $\sum\tau_{u,v}(v-u)\le1$, the series is absolutely convergent.

\begin{theorem}[Conservation with contact residuals]\label{v9:core:general}
For every such law, threshold, and positive coefficient,
\begin{equation}\label{v9:core:offcontact}
 (n+H^2)a+nc\rho-\frac{n\phi_z}{6}(\gamma+3J)+n^{3/2}L
 =\mathcal R_c+n^{3/2}(R_{\rm supp}-H\xi).
\end{equation}
At a fixed-target maximizing pair, both residuals vanish.
\end{theorem}
\begin{proof}
The gapwise identity in the proof of
Theorem~\ref{revision:contact-conservation} holds without contact.
Keeping its endpoint terms and the integral over $K$ adds exactly
$R_{\rm supp}$ to the right side of \eqref{v9:core:first}.
The zero-bias evaluation \eqref{v9:core:second} is unchanged. In the
first-moment calculation, $\E[Xh(X)]=\E[XQ(X)]-\xi$, so
\eqref{v9:core:firstmoment} becomes
\[
 M+\phi_z=-\frac{\gamma H\phi_z}{2n}+\frac{cK_1}{n}-\sqrt n\xi.
\]
The same pointwise identity for $q_z$ therefore gives
\[
 \sqrt n\mathcal A=c\rho+a-\gamma\phi_z/6+P_0/n
                    +\sqrt n\varepsilon/2+H\sqrt n\xi,
\]
with exactly the $P_0$ defined in the main proof. Equating the two
expressions and multiplying by $n$ gives
\eqref{v9:core:offcontact}. At contact, $\Delta=0$ on all of $K$,
so every term in both residuals vanishes, including the Lebesgue integral.
\end{proof}

\paragraph{Why the support integral is needed.}
Let $X$ be uniform on $[-\sqrt3,\sqrt3]$ and take $H=0$.
There are no support gaps, $F=T=1/2$, $\beta=\gamma=0$, and
$K_q=3\sqrt3/40$. Thus
\[
 R_{\rm supp}=\int_K\tau\Delta=cK_q/n^{3/2}>0.
\]
Dropping the support integral would give a false identity even though
all gap sums are empty. This example tests the extension without
contact; it does not contradict the contact theorem. The leakage
integral retains $[u,v)$ throughout: a right-endpoint atom belongs to
$h(v)$, and a left-endpoint atom has weight zero.
